\subsection{Differentiable Shape-Matching PBD Approach}
\label{sec:pbd-substrate}
We adopt a shape-matching PBD solver as our deformation substrate.
The choice is motivated by one architectural property specific to our
setting: its per-vertex participation weight $w_i$ enters every step
of the cluster solve as a primary input, providing the modulation
interface into which the continuous cut field of
\S\ref{sec:pf-cutting} is injected through a single assignment. We
review the substrate below in full to establish the notation and the
differentiability properties subsequent sections will rely on.


We work on a tetrahedral mesh
$\mathcal{M}=(\mathcal{V},\mathcal{E},\mathcal{F},\mathcal{T})$ and
identify each simulated particle with the mesh vertex it occupies, so a
single index $i\in\mathcal{V}$ labels both. The vertices are partitioned
into overlapping clusters $\{C_1,C_2,\ldots\}$, each a small group of
nearby vertices with a remembered rest shape. A cluster $C$ carries,
for every $i\in C$, a current position $\mathbf{p}_i\in\mathbb{R}^3$, a
rest position $\bar{\mathbf{p}}_i\in\mathbb{R}^3$, and a participation
weight $w_i\in\mathbb{R}_{\geq0}$. The rest positions
$\{\bar{\mathbf{p}}_i\}_{i\in C}$ are fixed at mesh initialization and
constitute the cluster's remembered shape; only $\mathbf{p}_i$ and
$w_i$ evolve during the simulation. The solve recovers the single
rigid rotation $\mathbf{R}_C\in\mathrm{SO}(3)$ that best aligns the
cluster's deformed configuration with its rest in a weighted
least-squares sense, in four closed-form steps.
 
\emph{(i)} The weighted centroids of the current and rest
configurations, $\boldsymbol{\mu}_C,\bar{\boldsymbol{\mu}}_C\in\mathbb{R}^3$,
\begin{equation}
  \boldsymbol{\mu}_C =
    \frac{\sum_{i \in C} w_i \mathbf{p}_i}{\sum_{i \in C} w_i},
  \qquad
  \bar{\boldsymbol{\mu}}_C =
    \frac{\sum_{i \in C} w_i \bar{\mathbf{p}}_i}{\sum_{i \in C} w_i},
  \label{eq:centroid}
\end{equation}
absorb the cluster's translation, leaving a pure rotation to find.

\emph{(ii)} The cross-covariance matrix
$\mathbf{A}_C\in\mathbb{R}^{3\times3}$,
\begin{equation}
  \mathbf{A}_C = \sum_{i \in C} w_i\,
     (\mathbf{p}_i - \boldsymbol{\mu}_C)
     (\bar{\mathbf{p}}_i - \bar{\boldsymbol{\mu}}_C)^\top,
  \label{eq:covariance}
\end{equation}
relates current deviations to rest deviations.

\emph{(iii)} The cluster rotation $\mathbf{R}_C\in\mathrm{SO}(3)$ is the
orthogonal factor of the polar decomposition
$\mathbf{A}_C=\mathbf{R}_C\mathbf{S}_C$, where $\mathbf{S}_C$ is
symmetric positive semidefinite and is not used further. We obtain
$\mathbf{R}_C$ with a short fixed iteration of quaternion
updates~\cite{muller2016}, whose parametrization guarantees
$\det\mathbf{R}_C=+1$ by construction and so avoids the $\mathrm{O}(3)$
reflection branch of SVD-based polar decomposition; the implicit
translation $\mathbf{t}_C=\boldsymbol{\mu}_C-\mathbf{R}_C
\bar{\boldsymbol{\mu}}_C\in\mathbb{R}^3$ then follows for free.

\emph{(iv)} Each vertex $i\in C$ is pulled toward a goal position
$\mathbf{g}_i^{(C)}$ and receives a correction $\Delta_i^{(C)}$,
\begin{equation}
  \mathbf{g}_i^{(C)} = \mathbf{R}_C \bar{\mathbf{p}}_i + \mathbf{t}_C,
  \qquad
  \Delta_i^{(C)} = w_i\,(\mathbf{g}_i^{(C)} - \mathbf{p}_i)
            \in \mathbb{R}^3,
  \label{eq:cluster-correction}
\end{equation}
the superscript $(C)$ marking that both quantities depend on the
cluster through $\mathbf{R}_C$ and $\boldsymbol{\mu}_C$.
 
What matters here is not the rotation fit itself, which is standard,
but what governs it. Every quantity in
(\ref{eq:centroid})--(\ref{eq:cluster-correction}) passes through the
participation weight $w_i$: it sets each vertex's pull on the
centroids, its influence on $\mathbf{R}_C$ through the covariance, and
the strength of its own correction. The weight is therefore the \emph{single
interface} through which a vertex's participation in its clusters can
be modulated, and the map from $\{w_i\}$ to the cluster correction is
differentiable. Away from degenerate clusters---collinear particles,
repeated singular values of $\mathbf{A}_C$, or vanishing total
weight---the operations form a smooth computation graph: the weighted
sums and outer product are bilinear, the quaternion iteration is
branch-free, and the goal map is linear. Two safeguards keep the
degenerate regime numerically benign: we floor the centroid divisor at
$\max(\sum_{i\in C} w_i,\,\varepsilon)$ against vanishing-mass clusters,
and add $\varepsilon\,\mathbf{I}_3$ to $\mathbf{A}_C$ before the polar
iteration, with $\mathbf{I}_3$ the $3\!\times\!3$ identity. Both
perturbations sit well below the simulation scale and, as
\S\ref{sec:gradient-verify} verifies, leave gradients unbiased in
finite-difference parity tests.
 
At each substep, $K$ such cluster solves run in parallel on the GPU in
a Jacobi pattern: every cluster reads a shared state, and the
overlapping per-vertex corrections are gathered into a fresh state for
the next iteration. Writing $c\in\{0,\dots,K\}$ for the Jacobi
iteration index and $\mathbf{p}^{(c)}$ for the corresponding position
iterate, the gathered update is applied with the global gain $k$, the
bulk stiffness:
\begin{equation}
  \mathbf{p}^{(c+1)}_i = \mathbf{p}^{(c)}_i +
    \frac{k}{|\mathcal{C}(i)|}\!\sum_{C \in \mathcal{C}(i)}
    \Delta_i^{(C)},
  \label{eq:apply}
\end{equation}
where $\mathcal{C}(i)$ is the set of clusters containing vertex $i$ and
each $\Delta_i^{(C)}$ is the correction (\ref{eq:cluster-correction})
evaluated at the current iterate $\mathbf{p}^{(c)}$. This substrate is
fully differentiable but leaves one discrete operation: cutting. The
naïve way to cut here---flipping $w_i$ to zero once a vertex is judged
severed---destroys exactly the differentiability in $w_i$ that the
cluster solve provides, as do the heavier remeshing and
vertex-duplication schemes of classical PBD cutting. 
